\chapter{Backend Design}
\section{RAG Pipeline Overview}

The backend implements a sequential RAG pipeline covering the full lifecycle from raw web content to streamed natural language responses. The pipeline is divided into an offline indexing phase and an online inference phase.

\paragraph{Offline indexing.}
Public web pages from the SUSTech official website are first collected by a crawling module that extracts raw HTML and persists it as structured records (\texttt{raw\_documents.jsonl}). A preprocessing stage then removes navigation elements, scripts, and redundant markup, yielding cleaned documents (\texttt{documents.cleaned.jsonl}). The cleaned text is subsequently segmented into semantically coherent chunks (\texttt{chunks.jsonl}) to improve retrieval granularity. Each chunk is encoded into a dense vector using a BGE embedding model \parencite{bge2024} and stored in a ChromaDB \parencite{chromadb} vector index.

\paragraph{Online inference.}
At query time, the user's question is encoded with the same embedding model and used to perform approximate nearest-neighbour search over the ChromaDB index, retrieving a candidate set of relevant chunks. A BGE reranker \parencite{bge2024} is then applied to re-score the candidates by cross-attention relevance, producing a refined context. The ranked context is assembled into a prompt and forwarded to the configured LLM backend—either llama.cpp \parencite{llamacpp} or vLLM \parencite{vllm2023}—which streams a Qwen3 \parencite{qwen2025} response token-by-token to the frontend via SSE.

\section{Pipeline Summary}

Table~\ref{tab:pipeline} summarises each stage of the pipeline, its input, output, and the responsible component.

\begin{table}[h]
\centering
\caption{RAG pipeline stages}
\label{tab:pipeline}
\begin{tabular}{llll}
\hline
\textbf{Stage} & \textbf{Input} & \textbf{Output} & \textbf{Component} \\
\hline
Crawling       & SUSTech web pages          & raw\_documents.jsonl      & Crawler \\
Preprocessing  & raw\_documents.jsonl       & documents.cleaned.jsonl   & Preprocessor \\
Chunking       & documents.cleaned.jsonl    & chunks.jsonl              & Chunker \\
Indexing       & chunks.jsonl               & ChromaDB vector index     & Indexer \\
Retrieval      & User query                 & Candidate chunks          & Retriever \\
Reranking      & Candidate chunks           & Ranked context            & Reranker \\
Generation     & Ranked context + query     & Streamed response         & LLM backend \\
\hline
\end{tabular}
\end{table}
